
\par\medskip\noindent\SGAInlineRefTarget{sga3:xxiv:proposition:5-5}\textbf{Proposition 5.5.---}\enspace
\emph{Let \(S\) be a scheme and \(G\) an adjoint semisimple
\(S\)-group.  There is a unique decomposition}
\[
  G\simeq\prod_{t\in\mathrm T}G_t,
\]
\emph{where \(G_t\) is an adjoint semisimple \(S\)-group, isotypic of
type \(t\).  Moreover, the preceding decomposition induces
isomorphisms}
\[
  \operatorname{Dyn}_t(G)\simeq\operatorname{Dyn}(G_t),
  \qquad
  \underline{\operatorname{Aut}}_{S\text{-}\mathrm{gr.}}(G)
  \simeq
  \prod_{t\in\mathrm T}
  \underline{\operatorname{Aut}}_{S\text{-}\mathrm{gr.}}(G_t).
\]

This was proved above when \(G\) is split.  In the general case one may
suppose that \(G\) has constant type \(\mathscr R\).  Using the
preceding decomposition of \(A_S(\mathscr R)\) and \SGARefLink{sga3:xxiv:1.17}{1.17}, one obtains
the desired decomposition of \(G\).  The other assertions then follow
by descent.

\par\medskip\noindent\SGAInlineRefTarget{sga3:xxiv:remark:5-6}\textbf{Remark 5.6.---}\enspace
More generally, if \(G\) and \(H\) are two adjoint semisimple groups,
there are canonical isomorphisms in the following commutative diagram:
\[
  \begin{tikzcd}[column sep=huge,row sep=large]
\underline{\operatorname{Isom}}_{S\text{-}\mathrm{gr.}}(G,H)
  \arrow[r,"\sim"] \arrow[d]
&
\displaystyle\prod_{t\in\mathrm T}
\underline{\operatorname{Isom}}_{S\text{-}\mathrm{gr.}}(G_t,H_t)
  \arrow[d]
\\
\underline{\operatorname{Isomext}}(G,H)
  \arrow[r,"\sim"] \arrow[d,"\wr"']
&
\displaystyle\prod_{t\in\mathrm T}
\underline{\operatorname{Isomext}}(G_t,H_t)
  \arrow[d,"\wr"]
\\
\underline{\operatorname{Isom}}_{\mathrm{Dyn}}
  \bigl(\operatorname{Dyn}(G),\operatorname{Dyn}(H)\bigr)
  \arrow[r,"\sim"]
&
\displaystyle\prod_{t\in\mathrm T}
\underline{\operatorname{Isom}}_{\mathrm{Dyn}}
  \bigl(\operatorname{Dyn}(G_t),\operatorname{Dyn}(H_t)\bigr).
\end{tikzcd}

\]

\par\medskip\noindent\SGAInlineRefTarget{sga3:xxiv:remark:5-7}\textbf{Remark 5.7.---}\enspace
One can give the following intrinsic characterization of \(G_t\),
which we state without proof: it is the largest reductive subgroup of
\(G\) that is invariant (and indeed characteristic) and isotypic of
type \(t\).

The preceding proposition reduces the study of adjoint semisimple
groups to that of isotypic adjoint semisimple groups.  We now undertake
this study.

\subsection*{\SGAInlineRefTarget{sga3:xxiv:subsection:5-8}5.8.}

Let \(\mathscr R\) be a pinned reduced adjoint simple root datum of type
\(t\), let \(I\) be a nonempty finite set, and let
\(\mathscr R^I\) be the product root datum of copies
\(\mathscr R_i\) of \(\mathscr R\), for \(i\in I\).  Denote by \(E\)
the group of automorphisms of \(\mathscr R\), identified with the group
of automorphisms of its Dynkin diagram \(D\).  The Dynkin diagram of
\(\mathscr R^I\) is identified with \(D^I\), and hence there is an
exact sequence
\[
  1\longrightarrow\operatorname{Aut}(D)^I
  \longrightarrow\operatorname{Aut}(D^I)
  \longrightarrow\operatorname{Aut}(I)
  \longrightarrow1,
\]
where \(\operatorname{Aut}(I)\) is the group of permutations of \(I\).
It follows that the canonical isomorphism
\[
  \operatorname{\acute Ep}_S(\mathscr R)^I
  \simeq
  \operatorname{\acute Ep}_S(\mathscr R^I)
\]
induces an exact sequence
\[
  1\longrightarrow A_S(\mathscr R)^I
  \longrightarrow A_S(\mathscr R^I)
  \longrightarrow\underline{\operatorname{Aut}}(I_S)
  \longrightarrow1.
\]
The last group is the constant \(S\)-group associated with
\(\operatorname{Aut}(I)\).  Notice also that \(I\) is canonically
identified with the set of connected components of \(D^I\).

Let \(G\) be an \(S\)-semisimple group of type \(\mathscr R^I\), thus
defining (cf.\ \SGARefLink{sga3:xxiv:1.17}{1.17}) a principal homogeneous bundle \(P\) under
\(A_S(\mathscr R^I)\).  The definitions by descent of
\(\operatorname{Dyn}(G)\) (\SGARefLink{sga3:xxiv:3.7}{3.7}) and of
\(\operatorname{Dyn}_0(G)\) (\SGARefLink{sga3:xxiv:subsection:5-2}{5.2}) show that the bundle associated with
\(P\) by
\[
  A_S(\mathscr R^I)\longrightarrow
  \underline{\operatorname{Aut}}(I_S)
\]
is precisely
\(\underline{\operatorname{Isom}}_S(I_S,\operatorname{Dyn}_0(G))\),
corresponding to the form \(\operatorname{Dyn}_0(G)\) of \(I_S\).
Using again the equivalence of categories \SGARefLink{sga3:xxiv:1.17}{1.17} and the preceding exact
sequence, one deduces formally that there is a reductive
\(\operatorname{Dyn}_0(G)\)-group \(G_0\) of type \(\mathscr R\), and
an \(S\)-isomorphism
\[
  G\simeq\prod_{\operatorname{Dyn}_0(G)/S}G_0.
\]

\par\medskip\noindent\SGAInlineRefTarget{sga3:xxiv:proposition:5-9}\textbf{Proposition 5.9.---}\enspace
\emph{Let \(S\) be a scheme and \(G\) an adjoint semisimple
\(S\)-group isotypic of type \(t\).  There exist an adjoint semisimple
\(\operatorname{Dyn}_0(G)\)-group \(G_0\), simple of type \(t\), and a
unique \(S\)-isomorphism}
\[
  G\simeq\prod_{\operatorname{Dyn}_0(G)/S}G_0.
\]
\emph{Moreover, this isomorphism induces an exact sequence}
\[
  1\longrightarrow
  \prod_{\operatorname{Dyn}_0(G)/S}
    \underline{\operatorname{Aut}}_{\mathrm{gr.}}(G_0)
  \longrightarrow
  \underline{\operatorname{Aut}}_{S\text{-}\mathrm{gr.}}(G)
  \longrightarrow
  \underline{\operatorname{Aut}}_S\bigl(\operatorname{Dyn}_0(G)\bigr)
  \longrightarrow1.
\]

One may plainly suppose that \(G\) has constant type \(n\,t\).  The
first assertion has already been proved (and uniqueness is clear).
The second then follows from the split case by descent.

Combining \SGARefLink{sga3:xxiv:proposition:5-5}{5.5} and \SGARefLink{sga3:xxiv:proposition:5-9}{5.9} gives:

\par\medskip\noindent\SGAInlineRefTarget{sga3:xxiv:proposition:5-10}\textbf{Proposition 5.10.---}\enspace
\emph{Let \(S\) be a scheme, let \(G\) be an adjoint semisimple
\(S\)-group, and put \(\mathfrak D=\operatorname{Dyn}(G)\).}

\emph{(i) There is a canonical decomposition}
\[
  G\simeq
  \prod_{t\in\mathrm T}
  \prod_{\mathfrak D_{0,t}/S}G_{0,t}
  \simeq
  \prod_{\operatorname{Dyn}_0(G)/S}G_0,
\]
\emph{where each \(G_{0,t}\) is a simple adjoint
\(\operatorname{Dyn}_{0,t}(G)\)-group of type \(t\); equivalently,
\(G_0\) is an adjoint semisimple
\(\operatorname{Dyn}_0(G)\)-group whose type at
\(x\in\operatorname{Dyn}_0(G)\) is \(a(x)\in\mathrm T\), where
\(a:\operatorname{Dyn}_0(G)\to\mathrm T_S\) was defined in \SGARefLink{sga3:xxiv:subsection:5-2}{5.2}.}

\emph{(ii) The preceding decompositions induce isomorphic exact
sequences, written vertically below.  Here
\(\underline{\operatorname{Aut}}_{S,a}(\mathfrak D_0)\) denotes the
scheme of automorphisms of \(\operatorname{Dyn}_0(G)\) commuting with
\(a:\operatorname{Dyn}_0(G)\to\mathrm T_S\):}
\[
  \begin{tikzcd}[column sep=large,row sep=large]
1 \arrow[d] & 1 \arrow[d] \\
\displaystyle
\prod_{t\in\mathrm T}\prod_{\mathfrak D_{0,t}/S}
  \underline{\operatorname{Aut}}_{\mathrm{gr.}}(G_{0,t})
  \arrow[r,"\sim"] \arrow[d]
&
\displaystyle
\prod_{\operatorname{Dyn}_0(G)/S}
  \underline{\operatorname{Aut}}_{\mathrm{gr.}}(G_0)
  \arrow[d]
\\
\underline{\operatorname{Aut}}_{S\text{-}\mathrm{gr.}}(G)
  \arrow[r,equal] \arrow[d]
&
\underline{\operatorname{Aut}}_{S\text{-}\mathrm{gr.}}(G)
  \arrow[d]
\\
\displaystyle\prod_{t\in\mathrm T}
  \underline{\operatorname{Aut}}_S
    \bigl(\operatorname{Dyn}_{0,t}(G)\bigr)
  \arrow[r,"\sim"] \arrow[d]
&
\underline{\operatorname{Aut}}_{S,a}
  \bigl(\operatorname{Dyn}_0(G)\bigr)
  \arrow[d]
\\
1 & 1
\end{tikzcd}

\]

\par\medskip\noindent\SGAInlineRefTarget{sga3:xxiv:corollary:5-11}\textbf{Corollary 5.11.---}\enspace
\emph{Under the preceding hypotheses, the following three categories
are equivalent:}
\begin{enumerate}
\item[(i)] the category of principal homogeneous bundles under \(G\);
\item[(ii)] the category of principal homogeneous bundles under \(G_0\);
\item[(iii)] the product, for \(t\in\mathrm T\), of the categories of
principal homogeneous bundles under the \(G_{0,t}\).
\end{enumerate}
This follows formally from the preceding decompositions and \SGARefLink{sga3:xxiv:proposition:8-4}{8.4}.

\par\medskip\noindent\SGAInlineRefTarget{sga3:xxiv:corollary:5-12}\textbf{Corollary 5.12.---}\enspace
\emph{There are canonical isomorphisms}
\[
  \underline{\operatorname{Tor}}(G)
  \simeq
  \prod_{t\in\mathrm T}
  \prod_{\mathfrak D_{0,t}/S}
    \underline{\operatorname{Tor}}(G_{0,t})
  \simeq
  \prod_{\operatorname{Dyn}_0(G)/S}
    \underline{\operatorname{Tor}}(G_0),
\]
\emph{and likewise with \(\underline{\operatorname{Tor}}\) replaced
by \(\underline{\operatorname{Bor}}\), respectively by
\(\underline{\operatorname{Kil}}\).}

This is immediate from the split case.

\par\medskip\noindent\SGAInlineRefTarget{sga3:xxiv:remark:5-13}\textbf{Remark 5.13.---}\enspace
The canonical morphism
\[
  \operatorname{Dyn}(G)\longrightarrow\operatorname{Dyn}_0(G)
\]
allows one to regard \(\operatorname{Dyn}(G)\) as a Dynkin scheme over
\(\operatorname{Dyn}_0(G)\); in fact, it is the Dynkin scheme
\(\operatorname{Dyn}(G_0)\) of \(G_0\).

Likewise, if \(T\subset B\) is a Killing couple of \(G\), corresponding
canonically to a Killing couple \(T_0\subset B_0\) of \(G_0\), then the
obstructions to the quasi-splitting of \(G\) and \(G_0\), which lie by
\SGARefLink{sga3:xxiv:3.9}{3.9} in
\[
  \operatorname{Pic}\bigl(\operatorname{Dyn}(G)\bigr)
  =
  \operatorname{Pic}\bigl(\operatorname{Dyn}(G_0)\bigr),
\]
coincide.  Hence:

\par\medskip\noindent\SGAInlineRefTarget{sga3:xxiv:corollary:5-14}\textbf{Corollary 5.14.---}\enspace
\emph{The following conditions are equivalent:}
\begin{enumerate}
\item[(i)] \(G\) is quasi-splittable;
\item[(ii)] \(G_0\) is quasi-splittable;
\item[(iii)] each \(G_{0,t}\), \(t\in\mathrm T\), is quasi-splittable.
\end{enumerate}

\section*{\SGAInlineRefTarget{sga3:xxiv:section:6}6. Automorphisms of Borel Subgroups of Reductive Groups}

\par\medskip\noindent\SGAInlineRefTarget{sga3:xxiv:lemma:6-1}\textbf{Lemma 6.1.---}\enspace
\emph{Let \(S\) be a scheme, let \((G,T,M,R)\) be a split \(S\)-group,
let \(\Delta\) be a system of simple roots of \(R\), and let \(B\) be
the corresponding Borel subgroup.  Then \(B^u\) is generated, as an
\textup{(fppf)} sheaf of groups, by the \(U_\alpha\),
\(\alpha\in\Delta\).}

Let \(H\) be the subgroup sheaf of \(B^u\) generated by the
\(U_\alpha\), \(\alpha\in\Delta\).  We prove
\(H\supset U_\beta\), for \(\beta\in R_+\), by induction on
\(\operatorname{ord}(\beta)=\operatorname{ord}_\Delta(\beta)>0\)
(cf.\ \SGARefLink{sga3:XXI:3.2.15}{Exp.~XXI, 3.2.15}).  The assertion holds by hypothesis when
\(\operatorname{ord}(\beta)=1\).  Suppose
\(\operatorname{ord}(\beta)>1\), and suppose
\(U_\gamma\subset H\) whenever
\(\operatorname{ord}(\gamma)<\operatorname{ord}(\beta)\).  There is an
\(\alpha\in\Delta\) such that \(\beta-\alpha\in R_+\)
(\SGARefLink{sga3:XXI:3.1.1}{Exp.~XXI, 3.1.1}).  Let \(p\) be the largest integer such that
\(\beta-p\alpha=\beta'\in R_+\).  Then
\(U_\alpha,U_{\beta'}\subset H\) and \(\beta'-\alpha\notin R\).
Thus it remains to prove:

\par\medskip\noindent\SGAInlineRefTarget{sga3:xxiv:lemma:6-2}\textbf{Lemma 6.2.---}\enspace
\emph{Keep the notation of \SGARefLink{sga3:XXIII:6.4}{Exp.~XXIII, 6.4}.  Suppose \(p=1\), that is,
\(\beta-\alpha\) is not a root.  If \(H\) is a subgroup sheaf of
\(B^u\) such that \(U_\alpha,U_\beta\subset H\), then
\(U_{i\alpha+j\beta}\subset H\) whenever
\(i\alpha+j\beta\in R\), \(i>0\), \(j>0\).}

Distinguish four cases according as \(q=0,1,2,3\).  In what follows,
\(x\) and \(y\) denote arbitrary sections of
\(\mathbb G_{a,S'}\), where \(S'\to S\).

For \(q=0\), there is nothing to prove.

For \(q=1\), one has
\[
 p_{\alpha+\beta}(\pm x)
 =p_\beta(-y)p_\alpha(-x)p_\beta(y)p_\alpha(x)\in H(S'),
\]
and hence \(U_{\alpha+\beta}\subset H\).

For \(q=2\), one similarly has
\[
 p_{\alpha+\beta}(\pm xy)
 p_{2\alpha+\beta}(\pm x^2y)
 =p_\beta(-y)p_\alpha(-x)p_\beta(y)p_\alpha(x)\in H(S').
\]
Thus, after changing some signs,
\[
 F(x,y)=p_{\alpha+\beta}(xy)p_{2\alpha+\beta}(x^2y)\in H(S').
\]
Since \(p_{\alpha+\beta}\) and \(p_{2\alpha+\beta}\) commute,
\[
 F(x,1)F(1,-x)=p_{2\alpha+\beta}(x^2-x).
\]
If \(a\in\mathbb G_a(S)\), the equation \(X^2-X=a\) defines a free
surjective extension of \(S\), namely
\(\operatorname{Spec}\mathcal O_S[X]/(X^2-X-a)\).  Consequently
\(p_{2\alpha+\beta}(a)\in H(S')\), so
\(U_{2\alpha+\beta}\subset H\), and then
\(U_{\alpha+\beta}\subset H\), since \(F(1,a)\in H(S')\).

For \(q=3\), one similarly has
\[
 F(x,y)=
 p_{\alpha+\beta}(xy)
 p_{2\alpha+\beta}(x^2y)
 p_{3\alpha+\beta}(x^3y)
 p_{3\alpha+2\beta}(x^3y^2)\in H(S'),
\]
and, since
\[
 p_{3\alpha+2\beta}(\pm x)
 =F(1,1)^{-1}p_\beta(-x)F(1,1)p_\beta(x)\in H(S'),
\]
one obtains \(U_{3\alpha+2\beta}\subset H\), and therefore
\[
 K(x,y)=
 p_{\alpha+\beta}(xy)
 p_{2\alpha+\beta}(x^2y)
 p_{3\alpha+\beta}(x^3y)\in H(S').
\]
Computing
\[
 K(x+y,1)K(-x,1)K(1,y)^{-1}
\]
modulo \(U_{3\alpha+2\beta}\), one finds
\[
 p_{2\alpha+\beta}(2x^2+y^2+2xy-y)
 p_{3\alpha+\beta}(y^3+3xy^2+3x^2y-y)\in H(S').
\]
If \(a\in\mathbb G_a(S)\), the ``equations''
\[
  x^2=-xy-y+1-a,\qquad
  y^2=3y-2+3a
\]
define a free surjective extension of \(S\).  One then has
\(y^3+3xy^2+3x^2y-y=0\), and the preceding expression gives
\(p_{2\alpha+\beta}(a)\in H(S')\).
Thus \(H\) contains \(U_{2\alpha+\beta}\) and
\(U_{3\alpha+2\beta}\), and
\[
 E(x,y)=p_{\alpha+\beta}(xy)p_{3\alpha+\beta}(x^3y)\in H(S').
\]
Since \(p_{\alpha+\beta}(xy)\) and \(p_{3\alpha+\beta}\) commute, the
preceding computation applies:
\[
 E(1,x)E(1,-x)=p_{3\alpha+\beta}(x^3-x).
\]
Hence \(U_{3\alpha+\beta}\subset H\), and then
\(U_{\alpha+\beta}\subset H\).

\par\medskip\noindent\SGAInlineRefTarget{sga3:xxiv:remark:6-2-1}\textbf{Remark 6.2.1.---}\enspace
The preceding proof shows that the hypothesis
``\(H\) contains \(U_\alpha\) and \(U_\beta\)'' could have been
replaced by: ``\(H\) contains \(U_\alpha\) or \(U_\beta\), and \(H\)
is invariant under \(\operatorname{int}(U_\alpha)\) and
\(\operatorname{int}(U_\beta)\).''

\par\medskip\noindent\SGAInlineRefTarget{sga3:xxiv:theorem:6-3}\textbf{Theorem 6.3.---}\enspace
\emph{Let \(S\) be a scheme, let \(G\) and \(G'\) be two semisimple
\(S\)-groups, and let \(B\) (respectively \(B'\)) be a Borel subgroup
of \(G\) (respectively \(G'\)).  Every isomorphism
\(B\xrightarrow{\sim}B'\) is induced by a unique isomorphism
\(G\xrightarrow{\sim}G'\).}

The assertion is local for the étale topology, so we may suppose that
\(G\) is split: \(G=(G,T,M,R)\), with \(B\) defined by the system
\(R_+\) of positive roots of \(R\).  The given isomorphism
\(u:B\xrightarrow{\sim}B'\) induces an isomorphism from \(T\) onto a
maximal torus \(T'\) of \(B'\), hence of \(G'\).  The given
isomorphism \(T\simeq D_S(M)\) gives an isomorphism
\(T'\simeq D_S(M)\) in which the elements of \(R_+\) become the
constant roots of \(B'\) relative to \(T'\), and hence the elements of
\(R\) become the constant roots of \(G'\) relative to \(T'\).  Since
\(G\) and \(G'\) are semisimple, the coroots are determined by duality
(\SGARefLink{sga3:XXI:1.2.5}{Exp.~XXI, 1.2.5}); this proves that \((T',M,R)\) is a splitting of
\(G'\) such that \(\mathscr R(G)=\mathscr R(G')\).

Applying \SGARefLink{sga3:XXIII:5.1}{Exp.~XXIII, 5.1} (the uniqueness theorem), one obtains a unique
isomorphism \(G\xrightarrow{\sim}G'\) agreeing with \(u\) on \(T\)
and on the \(U_\alpha\), \(\alpha\in\Delta\).  By \SGARefLink{sga3:xxiv:lemma:6-1}{6.1}, its restriction
to \(B\) is \(u\).

\par\medskip\noindent\SGAInlineRefTarget{sga3:xxiv:remark:6-3-1}\textbf{Remark 6.3.1.---}\enspace
(i) Using \SGARefLink{sga3:xxii:prop:4.1.9}{Exp.~XXII, 4.1.9} and arguing as in \emph{loc.\ cit.}, 4.2.12,
one may replace ``isomorphism'' in the theorem by ``isogeny''
(respectively, ``central isogeny'').

(ii) The theorem is false for reductive groups.  For example, take
\(G=G'=\mathrm{SL}_2\times\mathbb G_m\), identified with the group
\[
 \left\{
 \begin{pmatrix}
   a&b&0\\
   c&d&0\\
   0&0&h
 \end{pmatrix}
 \ \middle|\ ad-bc=1,\ h\text{ invertible}
 \right\}.
\]
